\documentclass[a4paper,11pt]{article}

\usepackage{revy}
\usepackage[utf8]{inputenc}
\usepackage[T1]{fontenc}
\usepackage[danish]{babel}
\usepackage{graphicx}


\revyname{MatematikRevyen}
\revyyear{2025}
% HUSK AT OPDATERE VERSIONSNUMMER
\version{0.1}
\eta{$3$ minutter}
\status{Færdig}

\title{Hulelignelsen}
\author{Frederikke '24, Louie '24, Ronni '25, Louise '23 og Julius '24}

\begin{document}
\maketitle

\begin{roles}
\role{M1} Matematiker
\role{M2} Matematiker
\role{F} Fysiker
\role{K} Kemiker
\end{roles}

\begin{props}
    \prop{Skillevæg}
    \prop{Bål} Behøver ikke nødvendigvis at være ægte ild
    \prop{Hævet scene} Evt. de store sorte appelboxes, eller nogle borde. 
    \prop{Lænker} til fysikeren og kemikeren.
    \prop{Skyggespil} Skygger på væggen under bandscenen. Kan evt. laves af pap og styres af ninjaer. Ellers kan det laves med en præcist fokuseret lampe (snak med \TeX nikken).
    \prop{Frugter} To appelsiner og et æble.
    \prop{Kasse} Lille kasse der kan indeholde katten.
    \prop{Kat} Evt. bare et tøjdyr.
    \prop{Heksagon} En perfekt konstrueret sekskant fra ideernes verden.
    \prop{F} F
\end{props}


\begin{sketch}

\scene{En skillevæg deler scenen op på midten. Til højre er en hævet scene hvorpå M1 og M2 står med et bål. M1 holder en appelsin, M2 holder et æble. På den anden side sidder F og K lænket fast så de blot kan se skyggerne på væggen kastet af lys fra bålet.}

\says{M1} Prøv at se det her trick, jeg har lært.

\scene{M1 drejer appelsinen på seks forskellige måder så den bliver til to.}

\says{M2} Wtf!? Hvad er det for et trylle trick. Hvordan har du snydt?

\says{M1} Er det ikke lige meget. Vil du have en appelsin?

\says{M2} Jo tak da. Så vil jeg ikke have det her æble.

\scene{smider æblet væk forbi bålet, så man kan se skyggen af det faldende æble}

\says{F} (fysikerlyde) OOOooooOOOOoooOOOOOohhhhhhh TYNGDEKRAFT!!!!

\says{M1} Kunne du også høre det? Det lød sygt dumt.

\says{M2} Jaja. Det bare fysikere. De lyder sådan. Prøv lige at se her:

\scene{M2 tager en kat og en kasse frem og holder katten foran bålet}

\says{F} (fysikerlyde) OOOooooOOOOoooOOOOOohhhhhhh EN KAT!

\scene{M2 holder kassen oppe}

\says{F} (fysikerlyde) OOOooooOOOOoooOOOOOohhhhhhh EN KASSE!

\scene{M2 holder katten og kassen oppe og putter katten ned i kassen}

\says{F} (fysikerlyde) OOOooooOOOOoooOOOOOohhhhhhh KVANTEMEKANIIIIIIK!!!!!

\says{M1} Hvordan lyder kemikere så?

\says{M2} Se her. Hvis jeg lige tager et heksagon frem.

\scene{Holder et heksagon oppe}

\says{K} (kemikerlyde) Whwhwhwhwhhwhwhwhwhwhwhrrrr BENZEEEEN!!

\says{M1} Hahaha jeg har en ide. Prøv lige at se hvis jeg giver dem en stamfunktion.

\scene{M1 holder et F oppe.}

\says{F} (fysikerlyde) OOOooooOOOOoooOOOOOohhhhhhh KRAAAAFT!!!
\says{K} (kemikerlyde) Whwhwhwhwhhwhwhwhwhwhwhrrrr FLOOUUUR!!!

\scene{F og K Kigger på hinanden anspændt stilhed}

\says{K} Det flour!

\says{F} Nej det kraft!

\says{M2} Nu får vi dem aldrig til at være stille igen.

\says{M1} Lad os stoppe mens legen er god.

\scene{Slår væggen ned og låser deres lænker op. F og K er skræmte}

\says{M1} \act{roligt som engle} Vil i med op i den abstrakte verden?

\says{M1 og M2} \act{Ud mod publikum} DEN VIRKELIGE VERDEN


\end{sketch}

\begin{figure}
    \centering
    \includegraphics[width=\linewidth]{IMG_20260425_153959.jpg}
\end{figure}
\end{document}
